%% 01_introduction.tex
\section{Introduction}
\label{sec:intro}

Specialty medications (biologics, long-acting injectables, and cyclic chemotherapy
regimens) now account for an estimated 51\% of US prescription drug spending despite
representing only a small fraction of claims~\citep{niu2024specialty}.  In Canada, biologics
alone constituted 29.6\% of public drug-plan spending from just 2.4\% of claims in
2022~\citep{cihi2023prescribed}.  Patient Support Programs (PSPs) provide the clinical
infrastructure that tracks adherence for these high-cost, administration-intensive
regimens~\citep{grundy2023psp}.  The questions PSP case managers ask every day (``When
is this patient's next scheduled injection?'', ``How many doses were administered in the
past 90 days?'', ``Has the patient missed any doses recently?'') map directly onto
canonical adherence indicators: Medication Possession Ratio (MPR), Proportion of Days
Covered (PDC), and persistence~\citep{lociganic2015using}.  Answering them requires
multi-step temporal reasoning over regimen structure that standard free-text search cannot
reliably perform~\citep{yao2024chemotimelines}.

Concurrently, LLM-powered ambient scribes have entered enterprise deployment at scale:
as of 2025--2026, systems from Abridge, Nuance DAX Copilot, and Ambience Healthcare
collectively serve millions of clinical encounters across academic health
systems~\citep{tierney2024ambient,rotenstein2026changes}.  These systems generate clinical
narratives as a by-product of documentation, creating a practical question: when the same
clinical facts are encoded in both a structured standard (FHIR) and an LLM-generated
narrative, which representation should a retrieval-augmented generation (RAG)
system~\citep{lewis2020retrieval} retrieve from?

Prior work has not answered this question because existing comparisons confound
representation with information content.  Studies on MIMIC-based predictive modelling
use independently authored structured and narrative sources for the same patient, so
performance differences reflect both representation and content coverage
simultaneously~\citep{moldwin2021empirical}.  DrugEHRQA pairs structured tables and
discharge notes for medication QA but explicitly acknowledges that ``information across the
two modalities is not strictly disjoint''~\citep{bardhan2022drugehrqa}.
FHIR-AgentBench~\citep{lee2025fhiragentbench} benchmarks agent architectures over
MIMIC-IV-FHIR but contains no narrative RAG arm.  LLMonFHIR~\citep{schmiedmayer2025llmonfhir}
evaluates only the structured function-calling variant.
CLEAR~\citep{lopez2025clinical} compares entity-aware vs.\ embedding RAG, but on
narrative notes, not FHIR resources.

\paragraph{This paper.}
We attempt to reduce the confound by constructing a paired dataset: 200 synthetic
patients each rendered as a FHIR~R4B bundle and as an LLM-generated narrative derived from
the same bundle, with 100\% recovery of eight prompted entity classes from the narrative
verified by a programmatic audit (\S\ref{sec:methods:dataset}).  Three retrieval systems
answer the same 13,800 questions from both representations.  We had originally described
the design as holding information content constant, so that accuracy differences would be
attributable to representation alone.  Reviewer scrutiny and a subsequent audit of our own
artefacts showed that this claim was too strong, in two directions: the narrative
contains a precomputed administration count that the FHIR serialisation does not state
explicitly, and the narrative omits most individual administration dates that the FHIR
bundle contains.  The narrative is also short enough (median 84 words) to fit in a single
retrieval chunk, so the narrative arm performs no retrieval in practice.  We therefore
present the comparison as a paired evaluation of two \emph{rendering pipelines} over a
common source, not as a clean isolation of representation.

\paragraph{What we found.}
The aggregate result is that the narrative arm scores highest (40.6\% vs.\ 35.3\% and
33.4\%).  Stratifying by whether the ground-truth answer is N/A, which is the case for
40.7\% of the question bank, shows that under the pre-specified scorer this lead comes
entirely from abstention: on questions whose correct answer is N/A the narrative arm is
far better at saying so, while on questions with a substantive answer all three arms score
15--16\% and are statistically indistinguishable under a patient-level cluster bootstrap.
A constant ``always N/A'' predictor matches the narrative arm overall.  Re-scoring only the
final answer, rather than the whole response including the reasoning trace, changes the
substantive picture: the narrative arm rises to 28.7\% while the structured arms stay near
17\%, and the residual lead sits in exactly the families whose answers the narrative
generation prompt writes into the text.  We therefore read the narrative advantage as a
combination of explicit statements of absence, summary leakage, and scorer behaviour, none
of which is a property of narrative as a format.  We also disclose three design flaws
found during the audit (a reference date that was never passed to the model, a
thinking-mode LLM under a tight output cap, and a scorer that reads the reasoning trace)
and quantify their reach.  We believe these cautionary findings are useful to the
clinical-RAG community precisely because the aggregate numbers looked like a clean
representation effect.

\paragraph{Contributions.}
\begin{enumerate}
\item \textbf{Paired-rendering benchmark.}  A released benchmark of 200 synthetic
  specialty-regimen patients rendered as FHIR~R4B bundles, LLM narratives, and templated
  narratives from a single source, with 13,800 programmatically verified questions across
  five PSP adherence-indicator families and three complexity tiers.  No LLM-as-judge is
  used anywhere.

\item \textbf{Resource-aware FHIR retriever.}  System~C adds a deterministic regex
  question router to a resource-type allowlist, one-hop reference-chain expansion, and
  temporal pre-filtering to standard dense retrieval, extending prior schema-aware
  retrieval approaches~\citep{jiang2023structgpt} to the FHIR R4B resource graph.

\item \textbf{Decomposition of the narrative advantage.}  We show that the
  narrative-over-structured gap on aggregate exact-match is an abstention effect under
  the pre-specified scorer, that it explains the apparent complexity gradient across
  tiers, and that the substantive-question lead which appears under final-answer scoring
  is confined to families whose answers the narrative summary states near-verbatim.

\item \textbf{Cluster-robust inference.}  All pairwise comparisons are re-estimated with
  a patient-level cluster bootstrap (200 clusters of 69 questions), and we report which
  effects narrow or lose significance relative to the question-level analysis.

\item \textbf{Audit of our own pipeline.}  We document three benchmark and pipeline
  defects (missing reference date, thinking-mode truncation, scorer reach into the
  reasoning trace), quantify their prevalence from the released raw outputs, and state
  which claims they invalidate.

\item \textbf{Open release.}  Code (Apache-2.0) and the 200-patient synthetic
  dataset (CC-BY-4.0) are released as a combined Zenodo deposit at
  \href{https://doi.org/10.5281/zenodo.20292792}{doi:10.5281/zenodo.20292792}.
\end{enumerate}
